\section{Controller Operations with MPC}
\label{sec:controller}

\subsection{Threshold Controller}

In the original Arca \texttt{st-contracts}, the controller was a single Ethereum address. If that key was compromised, the attacker could force-transfer any holder's securities. The Lux standard replaces the single controller with an MPC threshold controller.

Controller operations (forced transfers, redemptions, document updates) require a $t$-of-$n$ threshold signature from a controller quorum. The default configuration uses 2-of-3:

\begin{enumerate}[nosep]
  \item \textbf{Issuer share}: Held by the securities issuer.
  \item \textbf{Platform share}: Held by the ATS operator \cite{lux-ats-architecture}.
  \item \textbf{Regulatory share}: Held by the regulatory custodian or compliance officer.
\end{enumerate}

\begin{lstlisting}[language=Solidity,caption={MPC-controlled forced transfer.}]
contract SecurityToken is ISecurityToken {
    address public mpcController; // MPC threshold address

    modifier onlyController() {
        require(msg.sender == mpcController, "not controller");
        _;
    }

    function controllerTransfer(
        address from, address to, uint256 value,
        bytes calldata data, bytes calldata operatorData
    ) external override onlyController {
        // Controller can transfer regardless of restrictions
        // Used for: court orders, regulatory seizures, error correction
        _transfer(from, to, value);
        emit ControllerTransfer(msg.sender, from, to, value, data, operatorData);
    }

    function controllerRedeem(
        address holder, uint256 value,
        bytes calldata data, bytes calldata operatorData
    ) external override onlyController {
        _burn(holder, value);
        emit ControllerRedemption(msg.sender, holder, value, data, operatorData);
    }
}
\end{lstlisting}

The \texttt{mpcController} address is the public key derived from the distributed key generation ceremony. No single party ever possesses the full private key.

\subsection{Controller Action Types}

\begin{table}[H]
\centering
\caption{Controller actions and their legal bases.}
\label{tab:controller-actions}
\begin{tabular}{lp{6.5cm}}
\toprule
\textbf{Action} & \textbf{Legal Basis} \\
\midrule
Forced transfer & Court order, regulatory directive \\
Forced redemption & Issuer buyback, regulatory delisting \\
Account freeze & SAR investigation, OFAC designation \\
Document update & Prospectus amendment, offering circular update \\
Partition migration & Corporate action (merger, conversion) \\
Emergency pause & Material adverse event, security incident \\
\bottomrule
\end{tabular}
\end{table}

